% =====================================================================
% Appendix: Validation of the SimScale conjugate heat transfer model
% Requires: graphicx, subcaption, siunitx (v2 \SI syntax, as in main report)
% =====================================================================

\chapter{Conjugate Heat Transfer Model}
\label{app:cht_validation}

The SimScale conjugate heat transfer (CHT) model was checked in three stages
before its results were used: a mesh independence study to bound the
discretisation error, a check of iterative convergence, and a comparison of the
converged probe temperatures against the experimental measurements. The
acceptance criterion throughout was the $\pm\SI{2.2}{\celsius}$ tolerance of the
thermocouples. Numerical errors well inside this band can be neglected, so any
remaining disagreement with experiment can be attributed to the modelling
assumptions or to the measurements themselves.

\subsubsection{Flow Regime and Mesh Independence}

Turbulent cases were solved with the $k$--$\omega$ SST turbulence model, and a
laminar model was used for runs with an estimated $\mathrm{Re} \leq 2110$.
Because the two models resolve the near-wall region differently, mesh
independence was assessed separately for each flow regime and each coolant.
Each study was performed at the most demanding operating point in its regime,
the highest flow rate, where the boundary layers are thinnest and the near-wall
gradients steepest. A mesh that resolves this case is therefore conservative for
the lower flow rates in the same regime. Separate studies were needed for the
two coolants because their Prandtl numbers differ by more than two orders of
magnitude. Galinstan ($\mathrm{Pr} \sim 10^{-2}$) develops a thermal boundary
layer much thicker than its velocity boundary layer, whereas water
($\mathrm{Pr} \approx 5$--$7$) has a thermal boundary layer thinner than its
velocity boundary layer and so needs finer near-wall resolution to capture the
temperature field. The same sequence of five meshes, from \num{73e3} to
\num{1.5e6} cells, was used in every study.
% NOTE: your earlier paragraph said "three" meshes; the figure shows five.

Figure~\ref{fig:mesh_independence} shows the study for Galinstan at
\SI{1.4}{\milli\litre\per\second}. The inlet and outlet temperatures vary by
less than \SI{0.1}{\celsius} across the full range. The inlet is fixed by the
boundary condition and the outlet is set largely by the bulk energy balance, so
neither depends strongly on near-wall resolution. The resistor temperatures are
the sensitive quantities, because they depend on resolving conduction through
the cold plate and the temperature gradient at the wall. The two coarsest
meshes underpredict them by up to \SI{2.75}{\celsius} relative to the finest
mesh, which exceeds the thermocouple tolerance. From \num{427e3} cells upward,
every probe lies within \SI{0.25}{\celsius} of the \num{1.5e6}-cell result.
Convergence is not strictly monotonic: the \num{427e3}- and \num{655e3}-cell
meshes slightly overshoot the finest mesh, and the refinement ratio between
successive meshes is not constant. Richardson extrapolation and the grid
convergence index, which assume monotonic convergence at a fixed refinement
ratio, were therefore not applied. Instead, the \SI{0.25}{\celsius} residual
change was compared directly with the measurement tolerance, and it is about
one ninth of that tolerance. The \num{427e3}-cell mesh was adopted for the
Galinstan simulations because it is the coarsest mesh whose discretisation
error is negligible against the measurement uncertainty.
% CHECK: confirm 427k is the mesh actually used for the reported runs.

\begin{figure}[htbp]
    \centering
    \includegraphics[width=\textwidth]{mesh_independence.pdf}
    \caption{Mesh independence study for Galinstan at
        \SI{1.4}{\milli\litre\per\second}: (a) final probe temperatures and
        (b) change relative to the \num{1.5e6}-cell mesh. All probes change by
        less than \SI{0.25}{\celsius} between \num{427e3} and
        \num{1.5e6} cells.}
    \label{fig:mesh_independence}
\end{figure}

Figure~\ref{fig:mesh_independence_water} shows the corresponding study for
the turbulent water model at \SI{3.6}{\litre\per\minute}. The pattern is the
same: the inlet and outlet vary by less than \SI{0.03}{\celsius}, and the
resistor temperatures are underpredicted on the coarsest meshes, by up to
\SI{2.5}{\celsius} at \num{73e3} cells. The fine meshes settle more slowly than
they do for Galinstan, however. At \num{427e3} cells the resistors still differ
from the \num{1.5e6}-cell result by up to \SI{0.58}{\celsius}, and at
\num{655e3} cells by \SI{0.38}{\celsius}, compared with less than
\SI{0.25}{\celsius} for Galinstan on the same meshes. This is the behaviour the
Prandtl number argument predicts: the thinner thermal boundary layer in water
needs finer near-wall resolution before the wall temperature gradient, and
hence the resistor temperature, stops changing. Above \num{427e3} cells the
resistor temperatures fall monotonically with refinement, so the
\num{1.5e6}-cell result is itself probably still a few tenths of a degree from a
fully mesh-independent value. The \num{427e3}-cell mesh was retained for water.
Its discretisation error of about \SI{0.6}{\celsius} is roughly a quarter of the
thermocouple tolerance: acceptable for comparison with experiment, but not
negligible in the way it is for Galinstan.
% CHECK: confirm which mesh the water runs used. If 655k, replace the last two
% sentences: error about 0.4 C, roughly a sixth of the tolerance.
% NOTE: this is the turbulent water study. If a laminar water study (low flow)
% was also run, it needs its own figure or a sentence here.

\begin{figure}[htbp]
    \centering
    \figureorplaceholder[width=\textwidth]{mesh_independence_water.pdf}
    \caption{Mesh independence study for water at
        \SI{3.6}{\litre\per\minute} ($k$--$\omega$ SST): (a) final probe
        temperatures and (b) change relative to the \num{1.5e6}-cell mesh,
        plotted on the same scale as Figure~\ref{fig:mesh_independence}. The
        inlet and outlet differ by only \SI{0.26}{\celsius} and overlap in (a).
        The resistor temperatures change by up to \SI{0.58}{\celsius} between
        \num{427e3} and \num{1.5e6} cells.}
    \label{fig:mesh_independence_water}
\end{figure}
\issue{The measured full water flow is now \SI{3.5}{\litre\per\minute}
(\SI{58.3}{\milli\litre\per\second}) in the main text. This appendix still says
the water simulations ran at \SI{3.6}{\litre\per\minute}. If that is the flow
set in SimScale, keep it and say so; otherwise change it to
\SI{3.5}{\litre\per\minute}.}

\subsubsection{Iterative Convergence}

The model was solved as a steady-state problem, and convergence was judged by
monitoring the probe temperatures, since these are the quantities compared
against experiment (Figure~\ref{fig:simscale_full}). The inlet temperature
settles within about 150 iterations because it is fixed by the boundary
condition. The outlet and resistor temperatures relax far more slowly, decaying
over several thousand iterations after an initial overshoot. This behaviour is
typical of conjugate heat transfer: heat crosses the solid regions by
conduction alone, so the solid temperature field relaxes much more slowly than
the flow field. By 4000 iterations the outlet temperature had levelled off,
while the resistor temperatures were still falling by about \SI{0.3}{\celsius}
over the final 1000 iterations.
% CHECK: drift read from the plot; confirm from the exported probe data.
The fully converged resistor temperatures would therefore be slightly lower
than those reported. This drift is small relative to the measurement tolerance,
but it is not zero, and its direction matters for the comparison below: it would
move the right resistor towards the measured value and the left resistor
slightly further below it. The water model converges much faster
(Figure~\ref{fig:simscale_water}). After an initial transient, every probe is
within \SI{0.01}{\celsius} of its final value by about 1000 iterations, and at
every mesh density no probe changes by more than \SI{0.001}{\celsius} over the
final 1000 iterations. Iterative error is therefore negligible for water.

\begin{figure}[htbp]
    \centering
    \begin{subfigure}{\textwidth}
        \centering
        \includegraphics[width=0.9\textwidth]{simscale_convergence_full.pdf}
        \caption{Galinstan, \SI{1.4}{\milli\litre\per\second} (\SI{72}{\watt})}
        \label{fig:simscale_full}
    \end{subfigure}

    \vspace{0.6em}
    \begin{subfigure}{\textwidth}
        \centering
        \figureorplaceholder[width=0.9\textwidth]{simscale_convergence_water.pdf}
        \caption{Water, \SI{3.6}{\litre\per\minute} (\SI{72}{\watt}),
            \num{427e3} cells}
        \label{fig:simscale_water}
    \end{subfigure}
    \caption{Convergence of the SimScale conjugate heat transfer model against
        experiment. Solid lines show the simulated probe temperatures at each
        iteration; dotted lines show the measured values, with shaded bands for
        the $\pm\SI{2.2}{\celsius}$ thermocouple uncertainty. The measured
        resistor value is the mean of the two resistors. For Galinstan all
        probes lie within the measurement uncertainty at the final iteration;
        for water the resistors converge about \SI{2.4}{\celsius} below the
        measured value, just outside it.}
    \label{fig:simscale_convergence_both}
\end{figure}

\subsubsection{Comparison with Experiment}

For Galinstan, all four probes lie within the thermocouple tolerance of the
measured values at the final iteration (Figure~\ref{fig:simscale_full}). The inlet and outlet
agree to within about \SI{0.2}{\celsius}. The inlet agreement confirms that
the boundary condition was set correctly, rather than validating the model, and
the outlet agreement mainly confirms that the heat input and flow rate are
consistent with the energy balance. The resistor temperatures are the stronger
test, because they depend on the conduction path through the package and cold
plate and on the wall heat transfer coefficient the model predicts. The left
resistor is underpredicted by about \SI{0.5}{\celsius} and the right resistor
overpredicted by about \SI{1.4}{\celsius}. The model also predicts the right
resistor to be roughly \SI{2}{\celsius} warmer than the left, whereas the
measured temperatures are almost identical. This asymmetry lies within the
combined uncertainty, so the present instrumentation can neither confirm nor
rule it out.
\issue{The full-flow Galinstan data (\cref{fig:galinstan_transient}a) show the
right resistor at about \SI{60.1}{\celsius} and the left at about
\SI{58.6}{\celsius}, so the measured asymmetry is about \SI{1.5}{\celsius}, not
``almost identical''. The model's \SI{2}{\celsius} asymmetry is then close to
the measured one; reword this paragraph to match.}

The water model agrees less well (Figure~\ref{fig:simscale_water}). The
simulated inlet and outlet lie within \SI{0.15}{\celsius} of the measurements,
but this says little: the measured outlet rise is only \SI{0.4}{\celsius},
far below the sensor tolerance, so for water the outlet provides no useful
validation constraint. The resistors converge to \SI{46.9}{\celsius} and
\SI{46.8}{\celsius} against a measured mean of \SI{49.2}{\celsius}. They are
underpredicted by \SI{2.3}{\celsius} and \SI{2.4}{\celsius}, just outside the
tolerance band, and the resistor rise above the inlet is underpredicted by
about \SI{15}{\percent}. This gap is not a discretisation error. Refining the
mesh lowers the resistor temperatures further
(Figure~\ref{fig:mesh_independence_water}), so on the \num{1.5e6}-cell mesh the
underprediction grows to about \SI{2.9}{\celsius}. The discrepancy therefore
lies in the model itself. The cold plate, resistor package and their contact
are modelled identically for both coolants, and the Galinstan case shows no
comparable underprediction, which points towards the coolant side. The most
likely source is that the turbulent model overpredicts the wall heat transfer
coefficient for water. An error in the solid-side contact resistances that
happens to be masked in the Galinstan case cannot be excluded with only one
operating point per coolant.
% NOTE: measured values from Model Validation.xlsx ('FEM Model' sheet, water
% full flow): inlet 33.1, outlet 33.5, resistors 49.22 (mean). The raw log
% (CSV 17, 120-140 min) gives L 48.80, R 49.64 individually.

These results have clear limits. For Galinstan, agreement within
$\pm\SI{2.2}{\celsius}$ shows that the model is consistent with experiment at
the resolution of the measurements; it does not show that the model is
accurate to better than this. For water, the model is systematically optimistic,
predicting resistor temperatures about \SI{2}{\celsius} to \SI{3}{\celsius} too
low. The comparison also covers only one flow rate per coolant, so accuracy at
other operating points rests on the mesh study and the physical basis of the
model rather than on direct measurement. The model was therefore used where
these limits matter least: conclusions from the parametric studies are drawn
from trends and relative differences between configurations solved with the
same mesh and settings, which are less sensitive to systematic error than
absolute temperatures. The direction of the water bias also matters for the
coolant comparison. On average the model is slightly pessimistic for Galinstan and
optimistic for water, so any advantage the simulations show for Galinstan over
water is, if anything, understated.
% CHECK: confirm this matches how SimScale results are used in the main body.